% Options for packages loaded elsewhere
\PassOptionsToPackage{unicode}{hyperref}
\PassOptionsToPackage{hyphens}{url}
\PassOptionsToPackage{dvipsnames,svgnames,x11names}{xcolor}
\documentclass[
  10pt,
  fontset=fandol]{ctexart}
\usepackage{xcolor}
\usepackage[margin=16mm]{geometry}
\usepackage{amsmath,amssymb}
\setcounter{secnumdepth}{-\maxdimen} % remove section numbering
\usepackage{iftex}
\ifPDFTeX
  \usepackage[T1]{fontenc}
  \usepackage[utf8]{inputenc}
  \usepackage{textcomp} % provide euro and other symbols
\else % if luatex or xetex
  \usepackage{unicode-math} % this also loads fontspec
  \defaultfontfeatures{Scale=MatchLowercase}
  \defaultfontfeatures[\rmfamily]{Ligatures=TeX,Scale=1}
\fi
\usepackage{lmodern}
\ifPDFTeX\else
  % xetex/luatex font selection
\fi
% Use upquote if available, for straight quotes in verbatim environments
\IfFileExists{upquote.sty}{\usepackage{upquote}}{}
\IfFileExists{microtype.sty}{% use microtype if available
  \usepackage[]{microtype}
  \UseMicrotypeSet[protrusion]{basicmath} % disable protrusion for tt fonts
}{}
\makeatletter
\@ifundefined{KOMAClassName}{% if non-KOMA class
  \IfFileExists{parskip.sty}{%
    \usepackage{parskip}
  }{% else
    \setlength{\parindent}{0pt}
    \setlength{\parskip}{6pt plus 2pt minus 1pt}}
}{% if KOMA class
  \KOMAoptions{parskip=half}}
\makeatother
\setlength{\emergencystretch}{3em} % prevent overfull lines
\providecommand{\tightlist}{%
  \setlength{\itemsep}{0pt}\setlength{\parskip}{0pt}}
\usepackage{amsmath,amssymb,mathtools,needspace}
\xeCJKDeclareCharClass{CJK}{"2032,"0400->"04FF,"2160->"216B,"2460->"2473,"25A0,"25C7,"25CB,"25CF}
\IfFontExistsTF{Arial Unicode MS}{\xeCJKsetup{AutoFallBack=true}\setCJKfallbackfamilyfont{\CJKrmdefault}{Arial Unicode MS}}{}
\setcounter{secnumdepth}{0}
\setlength{\emergencystretch}{2em}
\allowdisplaybreaks
\usepackage{bookmark}
\IfFileExists{xurl.sty}{\usepackage{xurl}}{} % add URL line breaks if available
\urlstyle{same}
\hypersetup{
  pdftitle={Gauss 公式的稳定性},
  colorlinks=true,
  linkcolor={Maroon},
  filecolor={Maroon},
  citecolor={Blue},
  urlcolor={Blue},
  pdfcreator={LaTeX via pandoc}}

\title{Gauss 公式的稳定性}
\author{}
\date{}

\begin{document}
\maketitle

\subsubsection{4.4.4 Gauss 公式的稳定性}\label{gauss-ux516cux5f0fux7684ux7a33ux5b9aux6027}

对比 Newton-Cotes 公式，Gauss 公式不但是高精度的，而且是数值稳定的。Gauss 公式的稳定性之所以能够得到保证，是由于它的求积系数具有非负性。

\textbf{定理 4.6}　Gauss 公式（4.4.1）的求积系数 \(A_k\)（\(k=0,1,\cdots,n\)）全是正的。

\textbf{证明}　考察

\[
l_k(x)=\prod_{\substack{j=0\\j\ne k}}^{n}\frac{x-x_j}{x_k-x_j},
\]

它们是 \(n\) 次多项式，因而 \(l_k^2(x)\) 是 \(2n\) 次多项式，故 Gauss 公式（4.4.1）对于它能准确成立，即有

\[
\int_a^b l_k^2(x)\,\mathrm{d}x=\sum_{i=0}^{n}A_i l_k^2(x_i).
\]

注意到 \(l_k(x_i)=\delta_{ki}\)，上式右端实际上即等于 \(A_k\)，从而有 \(\displaystyle A_k=\int_a^b l_k^2(x)\,\mathrm{d}x>0\)。证毕。

现在讨论求积过程的稳定性问题。在用求积公式 \(\displaystyle I_n=\sum_{k=0}^{n}A_kf(x_k)\) 进行实际计算时，通常不一定能提供准确的数据 \(f_k=f(x_k)\)，而只是给出含有误差（譬如由于计

（本句续下页。）

\href{../../source-images/pdf-110.jpeg}{原页图像}

\subsubsection{4.4.4 Gauss 公式的稳定性（续）}\label{gauss-ux516cux5f0fux7684ux7a33ux5b9aux6027ux7eed}

算机字长的限制引进的舍入误差）的数据 \(f_k^*\)，故实际求得的积分值为

\[
I_n^*=\sum_{k=0}^{n}A_kf_k^*.
\]

人们自然关心，数据误差对积分值的影响能否加以控制呢？

由于 Gauss 公式的求积系数具有非负性，故

\[
|I_n^*-I_n|\leq\sum_{k=0}^{n}A_k|f_k^*-f_k|
\leq\left(\sum_{k=0}^{n}A_k\right)\max_{0\leq k\leq n}|f_k^*-f_k|.
\]

再利用式（4.1.4）的第一式，有

\[
|I_n^*-I_n|\leq(b-a)\max_{0\leq k\leq n}|f_k^*-f_k|,
\]

由此即可断定 Gauss 公式是稳定的。

\end{document}
